\subsection{两个格的相对不变量}

沿用第 2 至 5 节的记号和假设。设 V 是 K 上秩为 n 的向量空间，M 是 V 关于 A 的格。令 W 为外幂 \(\bigwedge^{n}V\)，它是 K 上秩为 1 的向量空间；令 \(M_{W}\) 表示 W 的格，该格由 \(M^{n}\) 在典范映射 \(V^{n}\to\bigwedge^{n}V\) 下的像生成（第 1 节，命题 3 (iii)；注意 \(M_{W}\) 不一定同构于 \(\bigwedge^{n}M\)（代数，第三章，§ 5，习题 9））。若 e 是 W 在 K 上的一个基，可以写 \(M_{W}=a.e\)，其中 a 是 A 的非零分式理想。\(\neq0\)

令 \(\mathbf{M}'\) 是 \(\mathbf{V}\) 的另一个格，并写作 \(\mathbf{M}_{\mathbf{W}}' = \mathfrak{a}'\) .e，其中 \(\mathfrak{a}'\) 是 \(\neq 0\)\(\mathbf{A}\) 的非零分式理想；除子 \(\operatorname{div}(\mathfrak{a}) - \operatorname{div}(\mathfrak{a}')\) 不依赖于所选的 \(e\)\(\mathbf{W}\) 的基 e，因为换基时 \(\mathfrak{a}\) 和 \(\mathfrak{a}'\) 同时乘以 \(\mathbf{K}^*\) 中的同一元素；我们写 \(\chi(\mathbf{M}, \mathbf{M}') = \operatorname{div}(\mathfrak{a}) - \operatorname{div}(\mathfrak{a}')\)，并称该除子为 \(\mathbf{M}'\) 关于 \(\mathbf{M}\) 的相对不变量。显然，若 \(\mathbf{M}, \mathbf{M}', \mathbf{M}''\) 是 \(\mathbf{V}\) 的三个格，则：

\[
\chi (\mathbf {M}, \mathbf {M} ^ {\prime}) + \chi (\mathbf {M} ^ {\prime}, \mathbf {M} ^ {\prime \prime}) + \chi (\mathbf {M} ^ {\prime \prime}, \mathbf {M}) = 0\tag{3}
\]

\[
\chi (\mathbf {M}, \mathbf {M} ^ {\prime}) + \chi (\mathbf {M} ^ {\prime}, \mathbf {M}) = 0\tag{4}.
\]

对所有 \(p \in P\)，由定义立即得到 \((M_{W})_p = (M_p)_{W}\)；此外，由于此时 \(M_p\) 是自由 \(A_p\)-模（因为 \(A_p\) 是主理想整环），\(M_p\) 在 \(A_p\) 上的一个基也是 V 在 K 上的一个基，因而 \(V\)\(K\)\((M_p)_W = \bigwedge^n (M_p)\)（第二章，§ 2，第 8 节），且分式理想 \(a_p = aA_p\) 是主理想整环。若令 \(a_p = p^{n_p} A_p\)、\(a' = p^{n'_{p}} A_p\)，则：

\[
\chi (\mathbf {M}, \mathbf {M} ^ {\prime}) = \sum_ {p \in P} (n _ {p} - n _ {p} ^ {\prime}). p,
\]

也可以写成：

\[
\chi (\mathbf {M}, \mathbf {M} ^ {\prime}) = \sum_ {p \in P} \chi (\mathbf {M} _ {p}, \mathbf {M} _ {p} ^ {\prime})\tag{5}
\]

将 \(\mathbf{D}(\mathbf{A}_{\mathfrak{p}})\) 识别为 \(\mathbf{D}(\mathbf{A})\) 中由 p 生成的子-Z-模。

命题 13. 设 M 是 V 的格，u 是 V 的一个 K-自同构。则：

\[
- \chi (\mathbf {M}, u (\mathbf {M})) = \operatorname{div} (\det (u)).\tag{6}
\]

对所有 \(p \in P\)，有 \(\bigwedge^{n} (u(M)_{p}) = \bigwedge^{n} (u(M_{p}))\)；若 \((e_{i})_{1 \leqslant i \leqslant n}\) 是 \(M_{p}\) 的一个基，则

\[
\bigwedge^ {n} \left(\mathrm{M} _ {p}\right) = \mathrm{A} _ {p}. e _ {1} \wedge e _ {2} \wedge \dots \wedge e _ {n},
\]

且 \(\bigwedge^{n}(u(\mathbf{M}_{\mathfrak{p}})) = \mathbf{A}_{\mathfrak{p}}.\det (u)e_{1}\wedge e_{2}\wedge \dots \wedge e_{n}\)，从而由公式（5）得到命题。

命题 14. 若 M、M' 是 V 的两个格，满足 M' ⊂ M，M/M' 是有限生成的挠 A-模，则

\[
\chi (\mathbf {M}, \mathbf {M} ^ {\prime}) = - \chi (\mathbf {M} / \mathbf {M} ^ {\prime}).\tag{7}
\]

显然 \(\mathbf{M} / \mathbf{M}' \subset \mathbf{V} / \mathbf{M}'\) 是有限生成的挠模；另一方面，对所有 \(\mathfrak{p} \in \mathbb{P}\)，我们知道（代数，第七章，§ 4，第 2 节，定理 1）存在  的基 \((e_i)_{1 \leqslant i \leqslant n}\) 和 \(\mathbf{M}_{\mathfrak{p}}\) 的基 \((e_i')_{1 \leqslant i \leqslant n}\)，使得当 \(\mathbf{M}_{\mathfrak{p}}'\) 时 \(e_i' = \pi^{\nu_i} e_i\)，其中 \(1 \leqslant i \leqslant n\)\(\nu_i \geqslant 0\)，而 \(\pi\) 是 \(\mathbf{A}_{\mathfrak{p}}\) 的一致化元。因此（沿用上面的记号）\(n_{\mathfrak{p}}' - n_{\mathfrak{p}} = \sum_{i=1}^{n} \nu_i\)；另一方面，\((\mathbf{M} / \mathbf{M}')_{\mathfrak{p}} = \mathbf{M}_{\mathfrak{p}} / \mathbf{M}_{\mathfrak{p}}'\) 同构于挠 \(\mathbf{A}_{\mathfrak{p}}\)-模 \(\bigoplus_{i=1}^{n} \mathbf{A}_{\mathfrak{p}} / p^{\nu_i} \mathbf{A}_{\mathfrak{p}}\)，因而其长度为 \(\sum_{i=1}^{n} \nu_i\)，结合（5）和第 5 节的定义 4，命题得证。

推论。设 \(L_{1}\)、\(L_{2}\) 是秩相同的两个自由 A-模，且 \(f: L_{1} \to L_{2}\) 是同态。令 U 为 f 关于 \(L_{1}\) 和 \(L_{2}\) 的基的矩阵。Coker(f) 为挠 A-模的充要条件是 \(\det(U) \neq 0\)，并且此时：

\[
\chi (\operatorname{Coker} (f)) = \operatorname{div} (\det (U)).\tag{8}
\]

\(L_{1}\) 和 \(L_{2}\) 可分别视为 \(V_{1} = L_{1} \otimes_{A} K\) 和 \(V_{2} = L_{2} \otimes_{A} K\) 中的格，f 延拓为 K-同态 \(f_{(K)}: V_{1} \to V_{2}\)。于是

\[
\left(\operatorname{Coker} (f)\right) _ {(\mathbf {K})} = \operatorname{Coker} \left(f _ {(\mathbf {K})}\right)
\]

而说 \(\operatorname{Coker}(f)\) 是挠 A-模，意味着 \(\operatorname{Coker}(f_{(\mathbf{K})}) = 0\)；现在，这等价于说 \(f_{(\mathbf{K})}\) 是满射，或者说 \(\det(U) \neq 0\)，从而得到第一断言。另一方面，可以写成 \(f(\mathbf{L}_1) = u(\mathbf{L}_2)\)，其中 u 是 \(u\)\(\mathbf{L}_2\) 的一个行列式为 \(\det(U)\) 的自同态；由于

\[
\operatorname{Coker} (f) = \mathrm{L} _ {2} / u (\mathrm{L} _ {2}),
\]

公式（8）由（7）和（6）得出。

例。如果 A = Z，则 A 的除子群可识别为有理数 0 的乘法群 \(Q_{+}^{*}\)。对于每个有限交换群 T，\textgreater\(\chi(T)\) 是 T 的阶；上述推论表明，群 \(\operatorname{Coker}(f)\) 的阶等于 \(\det(U)\) 的绝对值（参见代数，第七章，§4，第 7 节，定理 3 的推论 3）。
% semantic-fragments: legacy-input-seam
\relax{}
